% ---------------------------------------------------------------------
%  6 UŽDUOTIS — DI metodų efektyvumo palyginimas
%  REIKALINGA PREAMBULĖJE: tabularx, booktabs, siunitx, graphicx, float
%  Lentelės ir paveikslai GENERUOJAMI:
%    python -m src.eksperimentai.suvestine
%    python -m src.eksperimentai.pozymiu_svarba
%  Naujų eksperimentų šis skyrius nedaro: visi skaičiai imami iš
%  5 užduoties palikto vieno prėjimo per testavimo aibę.
% ---------------------------------------------------------------------

\subsection{Palyginimo pagrindas}
\label{sec:palyginimo_pagrindas}

\ref{sec:vertinimas} skyrius atsakė, kiek kiekvienas modelis pasiekia.
Šis skyrius atsako, kuris iš jų ir kuriomis sąlygomis tinkamiausias, o
tam reikia ne dar vieno metrikų rinkinio, o kompromisų: kiekvienas
pranašumas kuria nors ašimi perkamas nuolaida kita.

Palyginimas atliekamas ties suderintu klaidingų teigiamų biudžetu.
Priežastis nėra formali: ties didžiausios tikimybės tašku modelių
rikiuotė yra kitokia, o pats tas taškas šiame darbe atmestas, nes duoda
21--31~\% klaidingų teigiamų, t.~y. dvidešimt--trisdešimt
kartų daugiau, nei \ref{sec:atakos} skyriuje nustatytas eksploatacijos
biudžetas. Todėl
palyginimas ties argmax matuotų tašką, kurio sistema niekada nedirbtų.

Suvestinė turi dvi dalis. Trys prižiūrimi modeliai sprendžia aštuonių
kategorijų uždavinį, autokoderis --- dvejetainį, todėl bendro makro-F1
stulpelio visiems keturiems sudaryti negalima. Ši asimetrija numatyta dar
renkantis metodus (\ref{sec:di_metodai} skyrius) ir čia tik pateikiama
matomu pavidalu.

Skirtumai vertinami ta pačia taisykle kaip \ref{sec:vertinimas} skyriuje:
tikru laikomas tik toks, kuris viršija paleidimų sklaidą
\cite{dietterich1998tests}.


\subsection{Suvestinė}
\label{sec:suvestine}

\begin{table}[htbp]
\centering
\caption{Metodų palyginimas ties klaidingų teigiamų biudžetu.}
\label{tab:suvestine}
\lentele{suvestine.tex}
\end{table}

Prižiūrimų modelių rikiuotė ties operaciniu tašku yra XGBoost
(\num{0.663}), Random~Forest (\num{0.646}), MLP (\num{0.595}). Visų trijų
aptikimo lygis panašus --- 85--89~\% --- todėl skirtumą
daro ne tai, kiek atakų randama, o kaip tiksliai jos priskiriamos
kategorijoms ir kokia kaina biudžetas išlaikomas.

Du dalykai lentelėje svarbesni už rikiuotę. Pirma, \textbf{Random~Forest
biudžeto neišlaiko}: jo klaidingų teigiamų dalis testavimo aibėje yra
1,02~\%, nors slenkstis validacijos aibėje buvo parinktas
ties 0,94~\%. Antra, \textbf{modelių dydžiai skiriasi
tūkstantį kartų, o kokybė --- dešimtadaliu}: nuo 0,52~MB iki
638~MB makro-F1 pakyla nuo \num{0.595} iki \num{0.646}.
Būtent šis neproporcingumas, o ne rikiuotė, lemia rekomendaciją.

Autokoderis pateikiamas atskirai ir jo skaičiai su prižiūrimais
nelyginami tiesiogiai. Ties savo operaciniu tašku
(0,93~\% klaidingų teigiamų) jis pažymi kaip ataką
18,6~\% atakų srauto, o geriausiai atpažįstamoje
kategorijoje --- DDoS --- 30,8~\%. Tuo pačiu jo PR-AUC yra
\num{0.996}: rikiavimas geras, sprendimas ties reikalaujamu biudžetu ---
ne.


\subsection{Kompromisai}
\label{sec:kompromisai}

\begin{figure}[H]
  \centering
  \includegraphics[width=\textwidth]{kompromisai.pdf}
  \caption{Kairėje --- makro-F1 ties didžiausios tikimybės tašku ir ties
    klaidingų teigiamų biudžetu; skaičius virš stulpelių rodo santykinį
    kritimą. Dešinėje --- kokybė ties biudžetu prieš modelio dydį
    logaritmine skale.}
  \label{pav:kompromisai}
\end{figure}

\textbf{Kokybė prieš klaidingus teigiamus.} Biudžeto įvedimas kainuoja
visiems trims modeliams, bet nevienodai: XGBoost netenka
8,1~\% makro-F1, Random~Forest ---
11,3~\%, MLP --- 6,3~\%. Skirtumas
nedidelis absoliučiais dydžiais, bet jo pakanka rikiuotei apversti: ties
argmax pirmauja Random~Forest (\num{0.727} prieš \num{0.721}), ties
biudžetu --- XGBoost (\num{0.663} prieš \num{0.646}).

Priežastis yra tikimybių skiriamoji geba. Random~Forest tikimybė yra
medžių balsų dalis, todėl ties vienetu ji kinta šuoliais, o aukštas
slenkstis reikalauja būtent ten smulkios skalės; XGBoost grąžina
tolydžias reikšmes. Šis reiškinys išmatuotas trimis nepriklausomomis
sąlygomis --- validacijos aibėje, testavimo aibėje ir su nesuderintomis
hiperparametrų reikšmėmis --- todėl laikytinas metodų savybe, o ne vieno
matavimo ypatybe.

Iš to seka bendresnis metodinis teiginys: \emph{palyginimas ties
didžiausios tikimybės tašku ir palyginimas ties klaidingų teigiamų
biudžetu gali duoti skirtingą nugalėtoją}. Darbuose, kur skelbiamas tik
argmax rezultatas, šis skirtumas lieka nematomas.

\textbf{Kokybė prieš resursus.} Ši ašis skiria modelius stipriau nei
kokybė. Random~Forest yra 14 kartų didesnis už XGBoost
(638~MB prieš 45~MB) ir ties operaciniu tašku
prastesnis --- vadinasi, jis nepralaimi kompromise, o yra tiesiog
nurungtas: nė pagal vieną iš keturių atrankos kriterijų jis nėra
geresnis. Praktinė riba dar griežtesnė, nei rodo skaičius diske: modelis
į atmintį įkeliamas kelis kartus didesnis, todėl 1--8~GB
turinčiame kraštinio šliuzo įrenginyje netelpa iš principo.

Tikras kompromisas yra tarp XGBoost ir MLP. MLP yra
86 kartus mažesnis (0,52~MB) ir dešimt kartų
greitesnis, o kainuoja \num{0.068} makro-F1, t.~y.
10~\% santykinio kokybės. Įrenginiams, kuriems 45~MB
per daug --- mikrovaldiklių klasei, apie kurią kalba
\cite{mazinani2026constrained} --- tai reali alternatyva, o ne
nusileidimas.

Delsa nė vieno modelio neriboja: 3--30 mikrosekundžių
vienam įrašui yra trys eilės atsargos iki 20--50~ms
šliuzo biudžeto. Tai patvirtina \ref{sec:di_metodai} skyriuje iš
literatūros perimtą teiginį savais matavimais ir kartu paaiškina, kodėl
resursų kriterijų šiame uždavinyje realiai atstoja modelio dydis.

\textbf{Užduoties detalumas.} Detalumas yra ketvirtoji ašis, ir ji
vienintelė neduoda jokio kompromiso. Perėjus nuo aštuonių kategorijų prie
34 klasių, modelis padidėja 2,5 karto, mokymas pailgėja
3,6 karto, delsa auga 3,0 karto, o makro-F1 krenta
nuo \num{0.721} iki \num{0.646}. Smulkesnis detalumas kainuoja visomis
ašimis ir negrąžina nieko, todėl aštuonių kategorijų formuluotė yra ne
kompromisas, o geresnis pasirinkimas.

\textbf{Prižiūrimas prieš neprižiūrimą.} Autokoderio ir prižiūrimų
modelių skirtumas ties tuo pačiu klaidingų teigiamų biudžetu yra ne
laipsniškas: 18,6~\% prieš 85--89~\%
aptikimo. Kartu jo PR-AUC (\num{0.996}) rodo, kad anomalijos įvertis
atakas nuo gerybinio srauto \emph{rikiuoja} gerai. Abu teiginiai teisingi
ir neprieštarauja: rikiavimo kokybė virsta sprendimo kokybe tik tada, kai
egzistuoja slenkstis, ties kuriuo abu reikalavimai tenkinami vienu metu,
o gerybinio srauto atkūrimo paklaidos pasiskirstyme tokio taško nėra.


\subsection{Prognozės ir matavimas}
\label{sec:prognozes}

Metodai buvo parinkti \ref{sec:parinkimas} skyriuje pagal svertinę
sprendimų matricą, prieš atliekant eksperimentus. Prižiūrimų modelių
balai buvo XGBoost \num{4.70}, Random~Forest \num{3.55}, MLP \num{3.25};
matavimas ties operaciniu tašku davė \num{0.663}, \num{0.646} ir
\num{0.595} --- \textbf{tą pačią tvarką}. Tai prognozė, pateikta prieš
eksperimentą ir pasitvirtinusi, o ne paaiškinimas po fakto.

Sutapimas turi dvi išlygas, ir abi vertingesnės už patį sutapimą.

Pirma, \textbf{tvarka sutampa tik ties operaciniu tašku}. Ties argmax
matavimas duoda kitą pirmąją vietą, nors matrica yra ta pati. Vadinasi,
sprendimų matrica prognozavo ne šiaip metodo kokybę, o kokybę
sąlygomis, kurias pati matrica ir aprašo --- klaidingiems teigiamiems ten
skirta 30~\% svorio.

Antra, \textbf{vienas matricos įvertis buvo klaidingas}. Random~Forest už
resursus gavo 4 balus iš 5, o realiai jo modelis yra
638~MB ir į kraštinio šliuzo įrenginį netelpa. Klaida
sisteminė, ne atsitiktinė: balas remėsi algoritmo savybėmis, o modelio
dydis priklauso nuo mokymo aibės dydžio, kurio balų skalė neapima. Tas
pats pasirodė ir derinimo etape --- 400\,000 eilučių imtyje išmatuotas
196~MB dydis pilnoje aibėje virto 638~MB.

Lieka ir vienas nepatikrintas matricos teiginys: sprendimų medis surinko
\num{3.80} balo, t.~y. daugiau nei Random~Forest, bet į eksperimentą
nepateko, nes ketvertas buvo sudarytas ne vien pagal balus. Ar pigiausias
ir savaime paaiškinamas prižiūrimas metodas ties klaidingų teigiamų
biudžetu prilygtų gradientiniam stiprinimui, šis darbas atsakymo neturi.


\subsection{Nematytos atakos: hipotezės verdiktas}
\label{sec:zero_day}

Autokoderis į palyginimą įtrauktas ne dėl balo --- matricoje jis
pralaimėjo Isolation~Forest (\num{2.95} prieš \num{3.50}) --- o dėl
funkcinio reikalavimo aptikti atakas, kurių pavyzdžių mokymo aibėje
nebuvo. Toks pagrindimas yra tikrinamas teiginys, ir
\ref{sec:nematytos} poskyryje jis patikrintas.

Rezultatas neigiamas: prižiūrimas modelis, klasės niekada nematęs,
aptinka ją geriau dviem atvejais iš trijų, o trečiuoju abu metodai
lygūs. Palyginimui tai reiškia dvi išvadas.

Pirma, \textbf{neprižiūrimas metodas ketverte savo vietos nepateisino}.
Jis pralaimi ir pagrindinį uždavinį, ir tą, kuriam buvo įtrauktas. Iš
matricos seka, kad ketvirtoji vieta būtų buvusi naudingesnė
Isolation~Forest arba antram prižiūrimam metodui.

Antra --- ir tai svarbiau --- \textbf{apibendrinimo geba priklauso ne nuo
paradigmos, o nuo to, ar mokymo aibėje lieka gimininga klasė}. Kai
pašalinama viena klasė, o jos kategorija išlieka, prižiūrimo modelio
aptikimas nukrenta 0,0--1,5 procentinio punkto; kai
ištuštėja visa kategorija --- 21,3 punkto. Zero-day
klausimas todėl yra ne „prižiūrimas ar neprižiūrimas“, o „ar nauja ataka
patenka į jau žinomą šeimą“.

Teiginio ribos nurodytos \ref{sec:nematytos} poskyryje ir galioja
nepakitusios: tikrintos trys klasės, ir nė viena jų nebuvo naujo tipo
ataka, nepriklausanti jokiai iš aštuonių kategorijų.


\subsection{Kuo modelis remiasi}
\label{sec:pozymiu_svarba}

Interpretuojamumui atrankoje skirta 15~\% svorio, ir šis
kriterijus yra vienintelis, kurio prižiūrimų modelių palyginimas iki šiol
nelietė. Medžių ansamblis leidžia jį įvertinti be papildomo skaičiavimo:
vidutinis informacijos prieaugis apskaičiuojamas iš paties modelio, o ne
iš duomenų.

\begin{table}[htbp]
\centering
\caption{Dešimt daugiausiai informacijos prieaugio duodančių požymių.}
\label{tab:pozymiai}
\lentele{pozymiai.tex}
\end{table}

Pasiskirstymas labai netolygus: vienas požymis --- paketų skaičius
lange --- surenka 62,3~\% viso prieaugio, o penki
pirmieji kartu 80,4~\%. Tai atitinka imties sudėtį: DDoS,
DoS ir Mirai kategorijos kartu sudaro 70,8~\% eilučių, o
nuo gerybinio srauto jas skiria būtent intensyvumas. Tos pačios trys
kategorijos ties operaciniu tašku aptinkamos 100~\%. Kartu iš to seka
apribojimas --- modelis, kurio sprendimą dviem trečdaliais lemia vienas
intensyvumo požymis, žemo intensyvumo atakoms turi mažai atramos, ir
\ref{sec:nematytos} poskyryje aprašytas \texttt{DDOS-SLOWLORIS} atvejis
yra tiesioginė to iliustracija.

Antrą vietą užima protokolo kodas (6,0~\%). Tai
verta atskiro sakinio: rengiant požymių modulį šis stulpelis atrodė
perteklinis, nes rinkinyje jau yra atskiri TCP, UDP ir ICMP požymiai, ir
buvo tikrinamas kaip kandidatas šalinti. Patikra parodė, kad sutapimas
yra 70--94~\%, t.~y. koreliacija, o ne tapatybė, ir
stulpelis buvo paliktas. Svarbos lentelė šį sprendimą patvirtina
matavimu.

Šeši sąmoningai palikti mažos dispersijos požymiai kartu surenka
1,9~\% prieaugio. Jie naudojami visi, tad automatinis
filtras būtų pašalinęs informaciją, o ne triukšmą, tačiau įnašas mažas,
ir tai pasakytina taip pat aiškiai, kaip ir jų palikimo pagrindimas.

Metodo riba akivaizdi: informacijos prieaugis nerodo krypties --- jis
pasako, kuo modelis remiasi, o ne kokia požymio reikšmė reiškia ataką.
Sprendimų paaiškinimui analitikui to nepakanka, ir čia reikėtų
\cite{mohale2025xai} aprašytų priemonių, kurių skaičiavimo kaina yra
atskiro darbo dydžio. MLP ir autokoderis tokio pjūvio be papildomų
priemonių neduoda visai --- tai ir yra praktinė interpretuojamumo
kriterijaus išraiška.


\subsection{Rekomendacija}
\label{sec:rekomendacija}

Vieno atsakymo nėra, nes skiriasi diegimo sąlygos. Iš išmatuotų dydžių
seka keturios rekomendacijos.

\textbf{Kraštinio šliuzo įrenginiui, laikantis 1~\%
klaidingų teigiamų biudžeto, tinkamiausias yra XGBoost su sprendimo
slenksčiu.} Jis pasiekia makro-F1 \num{0.663}, aptinka
88,5~\% atakų prie 0,95~\% klaidingų
teigiamų, užima 45~MB ir vienam įrašui sunaudoja
29,6 mikrosekundės. Slenkstis turi būti parenkamas su
atsarga --- ties maždaug 80~\% biudžeto --- nes ties pačiu
kraštu parinktas taškas į naują duomenų aibę persikelia be atsargos, kaip
parodė \ref{sec:patikimumas} poskyrio patikra.

\textbf{Griežtai ribotos atminties įrenginiui tinka MLP.} Jis kainuoja
10~\% makro-F1, bet yra 86 kartus mažesnis ir
biudžeto laikosi patikimiausiai iš visų (0,64~\%).
Alternatyva --- mažesnė gradientinio stiprinimo konfigūracija: derinimo
paieškoje 200 medžių variantas atsiliko mažiau nei
1~\% ir buvo 7,4~MB, tačiau šis dydis matuotas
imtyje, o pilnoje mokymo aibėje jis paauga maždaug pusantro karto.

\textbf{Random~Forest šiam uždaviniui nerekomenduojamas.} Jis ties
operaciniu tašku prastesnis už XGBoost, vienintelis peržengia klaidingų
teigiamų biudžetą ir yra 14 kartų didesnis. Jo vertė šiame darbe --- ne
kaip sprendimo, o kaip atskaitos lygio, ties kuriuo matuojamas
gradientinio stiprinimo prieaugis.

\textbf{Autokoderis kaip vienintelis aptikimo metodas netinka.} Ties
reikalaujamu biudžetu jis aptinka mažiau nei penktadalį atakų srauto ir
nematytoms klasėms neaplenkia prižiūrimo modelio. Kaip papildoma
pakopa --- signalų rikiavimui arba antram vertinimui šalia prižiūrimo
modelio --- jis lieka svarstytinas, nes jo PR-AUC aukštas, o mokymas
trunka sekundes.

Rekomendacijų ribos nurodytinos kartu su jomis. Visi skaičiai gauti
viename duomenų rinkinyje, o kryžminis perkėlimas tarp rinkinių
literatūroje smarkiai krenta \cite{eren2026drift}, todėl rikiuotė
galioja šio pobūdžio srautui, o ne bet kuriam IoT tinklui. Zero-day
teiginys galioja naujai tos pačios šeimos klasei. Ir dvi svarbiausios
klaidų rūšys --- gerybinio srauto painiava su žvalgyba bei DoS ir DDoS
riba --- yra ne modelio, o duomenų savybė: joms atskirti reikalingų
krypties ir trukmės požymių šiame rinkinio leidime nėra, todėl jokia
metodo pakaita jų nepašalintų.
